\begin{center}
\textsc{Gauss} an \textsc{Bessel}.\quad Göttingen, 20.~November 1824.
\end{center}

\par\noindent\hspace*{\parindent}.\ .\ .\ .\ .\ .\quad Ich habe in diesem Jahre 12 Stationen besucht und bin mit den Dreiecken bis an die Weser (der entfernteste Punkt auf der Garlster Heide zwischen Osterholz und Vegesack) gekommen.
Die dritten Winkelpunkte liess ich anfangs immer zurück, und als ich nach der Mitte Augusts Bremen verliess, hatte ich noch 6 Plätze zu besuchen.
Mit allen ging es noch erträglich, aber bei dem letzten, dem Wilseder Berge, quälte mich das schlechte Wetter so, dass ich trotz einem dreiwöchentlichen Aufenthalt nicht ganz zu meiner Zufriedenheit fertig wurde; ich musste zuletzt einen Schlusstermin setzen und kam Ende Octobers nach Göttingen zurück.

\par\noindent\hspace*{\parindent}Ich behalte mir vor, Ihnen künftig ausführlicher über den Erfolg der Messungen, die grossen Schwierigkeiten, mit denen ich zu kämpfen gehabt (eine davon, das ewige Moorbrennen, wird Ihnen selbst noch im Gedächtniss sein), und manche interessante Phänomene, die sich dabei ergeben haben, zu schreiben.\quad.\ .\ .\ .\ .\ .

